\chapter{Исследовательская часть}

В данном разделе будут приведён пример работы программы, а также проведён сравнительный анализ алгоритмов на основе полученных данных.

\section{Технические характеристики}

Технические характеристики устройства, на котором выполнялись замеры времени, представлены далее.

\begin{enumerate}
	\item Процессор	AMD Ryzen 5 5500u ЦПУ 2592 МГц, ядер: 8, логических процессоров: 12.
	\item Оперативная память: 16 ГБ.
	\item Операционная система: Fedora 39~\cite{fedora}.
\end{enumerate}

При замерах времени ноутбук был включен в сеть электропитания. Запущены были только системные приложения.

\section{Демонстрация работы программы}

На рисунке \ref{img:working_example} представлена демонстрация работы разработанного приложения для алгоритмов поиска.

\includeimage
{working_example}
{f}
{H}
{1\textwidth}
{Пример работы программы}

\section{Число сравнений}

В результате проведения замеров по числу сравнений были получены столбчатые диаграммы~\ref{plt:dummy_search}--\ref{plt:bin_search_sorted}, при выполнении замеров элементы массива были отсортированы по возрастанию и принимали значения от 0 до 1024. Результат числа сравнений под индексом $-1$ является результатом поиска элемента, отсутствующего в рассматриваемом массиве.

\begin{figure}[H]
	\centering
	\includesvg[inkscapelatex=false, width=1\textwidth, height=1\textheight]{inc/img/DummySearch.svg}
	\caption{Зависимость количества сравнений от индекса искомого элемента в реализации алгоритма линейного поиска}
	\label{plt:dummy_search}
\end{figure}

\begin{figure}[H]
	\centering
	\includesvg[inkscapelatex=false, width=1\textwidth, height=1\textheight]{inc/img/BinSearch.svg}
	\caption{Зависимость количества сравнений от индекса искомого элемента в реализации алгоритма  бинарного поиска}
	\label{plt:bin_search}
\end{figure}

\begin{figure}[H]
	\centering
	\includesvg[inkscapelatex=false, width=1\textwidth, height=1\textheight]{inc/img/BinSearchSorted.svg}
	\caption{Алгоритм бинарного поиска с сортировкой по возрастанию количества сравнения элементов}
	\label{plt:bin_search_sorted}
\end{figure}

\section*{Вывод}
Из графиков~\ref{plt:dummy_search}--\ref{plt:bin_search_sorted}, можно сделать вывод, что максимальное число сравнений при использовании реализации бинарного поиска равно 10, а при использовании реализации линейного поиска~---~1024. Среднее количество сравнений для нахождения элемента в массиве алгоритмом бинарного поиска~---~8, линейного поиска~---~393. 

В ходе исследований получено, что алгоритм бинарного поиска лучше линейного поиска при большем размере исходного массива, поскольку требует в среднем меньше сравнений. При использовании бинарного поиска стоит учитывать необходимость упорядоченности элементов массива.


